% Step 9: Maximal-function selection, curvature-energy transport,
% and the Weitzenböck rotation constraint
%
% This note develops two mathematical constructions for the four-dimensional
% negative-branch collapse analysis.  Its fixed-cover regularity question and
% its homothetic-spectrum sentence were corrected in Step 10; see
% step10_multiscale_cover_and_diffuse_pontryagin.tex.  The maximal-function
% selection and Weitzenboeck energy calculation remain useful.

\section{Maximal-function good-point selection}

Fix a critical Ishida--LeBrun four-manifold $(M,g_i)$ with
\[
\delta_i = \|g_i^{-1}(\operatorname{Ric}_{g_i}+g_i)\|_\infty \to 0.
\]
The established near-saturation analysis provides:
\begin{itemize}
\item Total curvature energy bound:
$\int_M |\operatorname{Rm}_i|^2\,d\mu_i \le E_0(A)$
(from Chern--Gauss--Bonnet, signature, and the near-Einstein condition;
$A=\sum c_1^2(X_j)$ is the fixed topological parameter).
\item Calibrated deficit:
$\int_M F_i\,d\mu_i \to 0$
where $F_i = |\operatorname{Ric}_i+g_i|^2 + (s_i+8\sqrt6|W_i^+|)^2 + \mathcal{G}_i \ge 0$.
\item Total volume: $\operatorname{Vol}(M,g_i) \to V_\infty = \frac{32}{9}\pi^2 A > 0$.
\item Two-sided Ricci bound: $\operatorname{Ric}_i \ge -(1+\delta_i)\cdot 3 g_i \ge -4g_i$
for $i$ large, hence local volume doubling with constant $C_D = C_D(4,4)$.
\end{itemize}

\subsection{Hardy--Littlewood maximal function}
For a nonnegative $L^1$ function $f$ on $(M,g)$ with $\operatorname{Ric}\ge -(n-1)Kg$,
define
\[
\mathcal{M}_R f(p) = \sup_{0 < r \le R} \fint_{B_r(p)} f\,d\mu.
\]
The Ricci lower bound gives local volume doubling:
$\operatorname{Vol}(B_{2r}) \le C_D\operatorname{Vol}(B_r)$ for $r \le R$.
The standard Vitali-type covering argument then yields the weak-$(1,1)$ estimate:
\[
\operatorname{Vol}\bigl(\{\mathcal{M}_R f > K\}\bigr)
\le \frac{C_D}{K}\int_M f\,d\mu.
\]
This does not require noncollapse: only the doubling property, which follows from the
Ricci lower bound via Bishop--Gromov.

\subsection{Dual selection: curvature and deficit}
Apply this to both $f = |\operatorname{Rm}_i|^2$ and $f = F_i$ at scale $R = 1$.

\medskip\noindent
\textbf{Curvature-good set:}
\[
G_i^{\rm Rm}(K) = \bigl\{p \in M : \mathcal{M}_1 |\operatorname{Rm}_i|^2(p) \le K\bigr\}.
\]
Then $\operatorname{Vol}(M \setminus G_i^{\rm Rm}(K)) \le C_D E_0 / K$.

\medskip\noindent
\textbf{Deficit-good set:}
\[
G_i^F(\varepsilon) = \bigl\{p \in M : \mathcal{M}_1 F_i(p) \le \varepsilon\bigr\}.
\]
Then $\operatorname{Vol}(M \setminus G_i^F(\varepsilon)) \le C_D \int F_i / \varepsilon$.

\medskip\noindent
\textbf{Combined good set:} Choose $K = K_0 = 4 C_D E_0 / V_\infty$ (independent of $i$)
and $\varepsilon_i = \bigl(\int F_i\bigr)^{1/2} \to 0$. Then
\[
G_i = G_i^{\rm Rm}(K_0) \cap G_i^F(\varepsilon_i)
\]
satisfies
\[
\operatorname{Vol}(G_i) \ge V_\infty - \frac{V_\infty}{4}
- C_D \Bigl(\int F_i\Bigr)^{1/2} \to \frac{3V_\infty}{4}.
\]
For all large $i$, $\operatorname{Vol}(G_i) \ge V_\infty/2 > 0$.

\subsection{What the good set provides}
Every $p \in G_i$ satisfies, for all $0 < r \le 1$:
\begin{equation}\label{eq:good-point}
\fint_{B_r(p)} |\operatorname{Rm}_i|^2\,d\mu_i \le K_0,
\qquad
\fint_{B_r(p)} F_i\,d\mu_i \le \varepsilon_i \to 0.
\end{equation}
The constant $K_0$ depends on $A$ and $C_D$ but not on $i$ or $r$.
The deficit bound $\varepsilon_i$ tends to zero uniformly in $r$ and $p \in G_i$.
This is strictly stronger than having small energy at one radius.

\section{Transport to covering balls}

For $p \in G_i$, let $\widetilde{B}_1(\tilde p)$ be the local universal cover
of $B_1(p)$, equipped with the pulled-back metric.
The normal-cover mean-transfer lemma (Step~8, equation (T)) gives,
for $r \le 1/3$ and any nonneg.\ measurable $\phi$:
\[
\fint_{B_r(\tilde p)} \tilde\pi^*\phi\,d\tilde\mu
\le C_T \fint_{B_r(p)} \phi\,d\mu,
\qquad C_T = \frac{V_{-K}^{(4)}(1)}{V_{-K}^{(4)}(1/3)}.
\]
Applied to $\phi = |\operatorname{Rm}_i|^2$ and $\phi = F_i$:
\begin{equation}\label{eq:cover-good}
\fint_{B_r(\tilde p)} |\operatorname{Rm}_i|^2 \le C_T K_0,
\qquad
\fint_{B_r(\tilde p)} F_i \le C_T \varepsilon_i
\qquad\text{for all } r \le 1/3.
\end{equation}
The covering degree does not appear: it cancels between the multiplicity
numerator and the disjoint-deck-translate volume denominator.

\subsection{Absolute energy on the cover}
The absolute curvature energy on a covering ball is
\[
\int_{B_r(\tilde p)} |\operatorname{Rm}|^2\,d\tilde\mu
= \Bigl(\fint_{B_r(\tilde p)} |\operatorname{Rm}|^2\Bigr)
\cdot \operatorname{Vol}(B_r(\tilde p))
\le C_T K_0 \cdot \operatorname{Vol}(B_r(\tilde p)).
\]
If the cover has a volume upper bound $\operatorname{Vol}(B_r(\tilde p))\le C_Vr^4$,
then the absolute energy is at most $C_TK_0C_Vr^4$, hence small for small fixed
$r$.  Noncollapse supplies the separate lower volume ratio needed for a
harmonic chart.  If the cover remains collapsed, small absolute energy alone
does not provide such a chart.

\section{The regularity question}

\subsection{Precise statement}
\textbf{Question.} Does the following hold?
For every $K > 0$, there exists $r_*(K) > 0$ such that if $(B_1(\tilde p), \tilde g)$
satisfies $\operatorname{Ric} \ge -3\tilde g$ and
\[
\fint_{B_r(\tilde p)} |\operatorname{Rm}|^2 \le K
\quad\text{for all } 0 < r \le 1/3,
\]
then $\operatorname{Vol}(B_{r_*}(\tilde p)) \ge c(K) r_*^4$?

If yes, combined with \eqref{eq:cover-good}, this would immediately
give a harmonic chart at scale $r_*$, and the small deficit
$F_i \to 0$ would produce the forbidden calibrated Einstein germ.
The calibrated-germ contradiction survives at scale $r_*$ because
the rescaled Einstein constant is $-r_*^2 \ne 0$.

If no, the cover itself is collapsed despite controlled normalized
curvature energy at every scale. A countermodel would reveal
exactly what additional curvature or topological structure is needed.

\subsection{What is known}
\begin{itemize}
\item Anderson's $\varepsilon$-regularity: small absolute $\int|Rm|^2 < \varepsilon_0$
plus volume lower bound gives a harmonic chart. This requires noncollapse as input.
\item Cheeger--Naber: noncollapse + bounded $|\operatorname{Ric}|$ gives
codim-4 regularity (Thm~1.4) and, in 4D, finite diffeo types (Thm~1.12).
These also require noncollapse as input.
\item Naber--Zhang: under specific GH-model, rank, and model-noncollapsing-radius
conditions, the local universal cover \emph{is} noncollapsed with controlled
harmonic radius (Prop~6.4). The question is whether our controlled maximal
curvature function provides these conditions.
\item The 4D conformal invariance $\int|Rm|^2 = \text{(scale-invariant)}$
is consistent but does not by itself produce noncollapse.
\end{itemize}

\subsection{Choosing a small but fixed scale}
The calibrated-germ contradiction requires Ric $\approx -g$, not Ric $\approx 0$.
After rescaling $B_r \to B_1$ by $\tilde g = r^{-2}g$:
\[
\operatorname{Ric}_{\tilde g} = \operatorname{Ric}_g = -g + O(\delta) = -r^2\tilde g + O(\delta).
\]
For $r$ fixed and $\delta_i \to 0$, the rescaled limit satisfies
$\operatorname{Ric} = -r^2 g$, hence $s = -4r^2 \ne 0$.
As a $(1,3)$ tensor $W$ is conformally invariant under constant rescaling,
but as the curvature endomorphism $W^+:\Lambda^+\to\Lambda^+$ its spectrum
scales by $r^2$.  Thus the calibrated spectrum becomes
$(-r^2/6,r^2/12,r^2/12)$.  The projector obstruction and Einstein--Weyl
identity remain contradictory for every fixed $r>0$; only $r\to0$ sends the
Einstein scale and the spectral values to zero.

\section{Weitzenb\"ock gradient energy of harmonic self-dual forms}

\subsection{The formula}
On a 4D Einstein manifold with $\operatorname{Ric} = \lambda g$ ($s = 4\lambda$),
the Weitzenb\"ock formula for 2-forms (Besse, Theorem 1.143/1.155) is
\[
\Delta_d\omega = \nabla^*\nabla\omega + q(R)(\omega),
\]
where, for Einstein metrics:
\[
q(R)(\omega)_{ij}
= \underbrace{R_{ik}\omega_{kj} + R_{jk}\omega_{ik}}_{\text{Ricci terms: } 2\lambda\omega_{ij}}
- 2\underbrace{R_{ikjl}\omega_{kl}}_{\widetilde{R}(\omega)_{ij}}.
\]
The endomorphism $\widetilde{R}$ uses the index arrangement
$R_{ikjl}$, \textbf{not} the curvature-operator arrangement $R_{ijkl}$.
Its eigenvalues on $\Lambda^2_+$ are \emph{not} the usual $W^+ + s/12$.

For a harmonic form $H$ ($\Delta_d H = 0$) in the $\mu$-eigenspace
of $\widetilde{R}|_{\Lambda^2_+}$:
\[
\nabla^*\nabla H = -(2\lambda - 2\mu)H,
\qquad
\int|\nabla H|^2 = -(2\lambda - 2\mu)\int|H|^2.
\]

\subsection{Numerical verification at the calibrated spectrum}

A direct index computation (script: \texttt{verify\_weitzenbock\_clean.py})
on an explicit Einstein metric with $\lambda = -1$ and
$\operatorname{Spec}(W^+) = (-1/6, 1/12, 1/12)$ gives:

\begin{center}
\begin{tabular}{lccc}
\hline
Eigenform & $W^+$ eigenvalue & $\widetilde{R}$ eigenvalue $\mu$
& $\int|\nabla H|^2 / \int|H|^2$ \\
\hline
$\sigma_1$ (simple) & $-1/6$ & $-1/2$ & $\mathbf{1}$ \\
$\sigma_2$ (double) & $1/12$ & $-1/4$ & $\mathbf{3/2}$ \\
$\sigma_3$ (double) & $1/12$ & $-1/4$ & $\mathbf{3/2}$ \\
\hline
\end{tabular}
\end{center}

\medskip\noindent
\textbf{Consistency checks:}
\begin{itemize}
\item K\"ahler--Einstein with $\lambda=-1$: parallel K\"ahler form $\omega$
has $\widetilde{R}(\omega) = -1\cdot\omega$, giving
$q(R)(\omega) = 2(-1)\omega - 2(-1)\omega = 0$. Consistent with $\nabla\omega = 0$. $\checkmark$
\item K3 surface ($\lambda = 0$, $W^+ = 0$): $\widetilde{R} = 0$,
$q(R) = 0$, harmonic = parallel. $\checkmark$
\item Round $S^4$ ($\lambda = 3$, $W^+ = 0$): $\widetilde{R}$ eigenvalues
make $-(2\lambda-2\mu) < 0$, so no harmonic SD forms.
Consistent with $b^+_2(S^4) = 0$. $\checkmark$
\end{itemize}

\subsection{Key conclusion}
\textbf{At the calibrated Einstein spectrum, all harmonic self-dual 2-forms
have strictly positive gradient energy.} The gradient energy is at least
$\int|H|^2$ (for the simple eigenform) and $\frac32\int|H|^2$
(for the double eigenspace). In particular:

\medskip\noindent
\textit{The near-harmonic self-dual forms on the near-saturation sequence
are not nearly parallel. Their covariant derivative has $L^2$ norm
of the same order as their amplitude.}

\medskip\noindent
This corrects any implicit assumption from the Codex decomposition
that the forms might be approximately parallel. The Codex analysis
correctly established near-constant amplitude, approximate harmonicity,
and controlled eigenline alignment, but ``nearly parallel'' does not follow
and is in fact false.

\section{Implications for collapse geometry}

\subsection{Controlled rotation of the eigenline}
The simple eigenform $\sigma_1$ of $W^+$ (eigenvalue $-1/6$) defines
a section of the projectivized SD bundle $\mathbb{P}(\Lambda^2_+)$.
The Weitzenb\"ock identity shows this section has controlled twist:
$\int|\nabla\sigma_1|^2 = \int|\sigma_1|^2$ at the calibrated limit,
and $\int|\nabla\sigma_1|^2 = (1 + O(\delta_i))\int|\sigma_1|^2$
for the near-Einstein sequence.

Since $|\sigma_1| \approx \mathrm{const}$ (from the amplitude control):
\[
\fint_M |\nabla\sigma_1|^2\,d\mu \approx |\sigma_1|^2_\infty.
\]
The pointwise covariant derivative is, on average, comparable to the
amplitude. The eigenline direction rotates at an $O(1)$ rate.

\subsection{Rotation is a base phenomenon in fibered collapse}
In a regular fibered collapse with fiber $F$ of size $\ell \to 0$,
decompose $\nabla\sigma_1 = \nabla_B\sigma_1 + \nabla_V\sigma_1$.
The fiber holonomy of the $\Lambda^2_+$ bundle around a closed
fiber loop of length $\sim\ell$ is $O(|\operatorname{Rm}|\cdot\ell^2)$.
For bounded curvature, this is $O(\ell^2)$, giving
$|\nabla_V\sigma_1| \sim O(\ell)\cdot|\sigma_1|$ and
$\int|\nabla_V\sigma_1|^2 = O(\ell^2)\int|\sigma_1|^2 \to 0$.

Therefore the horizontal gradient must carry the full energy:
\[
\int|\nabla_B\sigma_1|^2 \approx \int|\sigma_1|^2.
\]
The $O(1)$ rotation rate lives in the base directions, not the collapsing
fibers. This means the base geometry (a 1-, 2-, or 3-dimensional space)
must support a section of the SD eigenline bundle with substantial
covariant derivative energy.

\subsection{No contradiction from gradient energy alone}
The positive gradient energy does not by itself exclude collapse.
A collapsed metric can accommodate rotating self-dual eigenlines.
However, on any \emph{noncollapsed regular patch} where the metric
converges to a smooth Einstein limit, the full Einstein equations
force the projector obstruction: constant calibrated $W^+$ spectrum
requires $\nabla W^+ = 0$, which requires $|\nabla W^+|^2 = 0$,
contradicting the Einstein--Weyl identity's requirement
$|\nabla W^+|^2 = 1/16$.

The gradient energy computation adds precision to the geometric picture:
\textit{the collapsed sequence must accommodate a controlled rotating
eigenline while avoiding every regular patch on which the full Einstein
equations would make that rotation impossible.}

\section{Combined framework and the remaining question}

\subsection{What is established}
On the fixed critical Ishida--LeBrun manifold with $\delta_i\to0$:
\begin{enumerate}
\item A combined good set $G_i$ of volume $\ge V_\infty/2$ exists
where both $\mathcal{M}_1|\operatorname{Rm}|^2 \le K_0$
and $\mathcal{M}_1 F_i \le \varepsilon_i\to0$.
\item On the local universal cover of any good-point ball,
both normalized means transport without sheet-number loss.
\item If the cover at any fixed scale $r_* > 0$ is noncollapsed,
the harmonic chart, small deficit, and projector obstruction
combine to give a contradiction.
\item All harmonic SD forms have $O(1)$ gradient energy per unit
amplitude at the calibrated spectrum, forcing eigenline rotation
at a controlled rate. This rotation lives in the base directions.
\end{enumerate}

\subsection{What remains}
The corrected question is whether the simultaneous curvature and calibrated
deficit bounds permit one to choose a buffered downstairs domain, before
taking its normal cover, at some scale bounded below in terms of $K$, so that
the cover is noncollapsed.  A cover fixed from an oversized domain is not the
right invariant object.

Stated precisely: does $\fint_{B_r(\tilde p)}|\operatorname{Rm}|^2 \le K$
for all $r \le 1/3$, combined with $\operatorname{Ric}\ge -3g$,
imply $\operatorname{Vol}(B_{r_*}(\tilde p)) \ge c(K)r_*^4$
for some $r_*(K) > 0$?

If yes: the four-dimensional negative-branch fixed-topology pinching
gap is proved.

If no: the failure mode---the cover itself being collapsed despite
controlled mean curvature at every scale---becomes a precise
geometric object to study. The eigenline rotation constraint
then serves as an additional structural input.

\subsection{Files}
\begin{itemize}
\item Verification script: \texttt{workspace/scripts/verify\_weitzenbock\_clean.py}
\item Previous curvature-energy script: \texttt{workspace/scripts/verify\_weitzenbock\_gradient\_energy.py}
\item Step 8 cover-mean transport: \texttt{workspace/notes/step8\_cover\_mean\_transport\_and\_model\_scope\_corrections.md}
\item Step 6 projector obstruction: \texttt{workspace/notes/step6\_local\_projector\_obstruction\_and\_collapse\_reduction.tex}
\end{itemize}
